\subsection{Redshift Binning and Photo-z Effects}\label{subsec:obs_redshift_binning}

This section addresses how redshift information enters observables through photometric redshifts and redshift bins.

\subsubsection{Redshift distribution}

For a photometric survey, sources are divided into bins $i, j$ with distribution $\frac{dn_i}{dz}(z)$.

The lensing kernel for bin $i$ is:

\begin{equation}
W^i(\chi) = \int dz \, \frac{dn_i}{dz}(z) \, W(\chi(z))
\end{equation}

where $W(\chi)$ is the lensing efficiency per unit distance.

\subsubsection{Cross-spectrum between bins}

Angular power spectra between different redshift bins are:

\begin{equation}
C_\ell^{ij} = \int_0^\infty \frac{dk}{2\pi^2} k^2 P(k) \left[\int d\chi \, \phi_i(\chi) j_\ell(k\chi)\right] \left[\int d\chi' \, \phi_j(\chi') j_\ell(k\chi')\right]
\end{equation}

If bins overlap in redshift, $C_\ell^{ij} \neq 0$ even for $i \neq j$ (correlation from common sources).

\subsubsection{Photo-z uncertainties}

Photometric redshifts have errors $\sigma_z \sim 0.01--0.05$ (relative). Effects:
- Broadens effective $\phi_i(\chi)$
- Causes bin-to-bin cross-contamination
- Biases $C_\ell$ measurements at level $\sim \sigma_z^2$

Must include photo-z covariance in analysis.

See \cref{subsec:obs_cross_spectra} for more on bin combinations.
